% Part: lambda-calculus
% Chapter: introduction
% Section: fixed-point-combinator

\documentclass[../../../include/open-logic-section]{subfiles}

\begin{document}

\olfileid{lam}{int}{fix}
\olsection{مؤلفات النقطة الثابتة}

لتكن لدينا عبارة لامبدا~$g$، ولنطلب حدًا آخر~$k$ يكون فيه~$k$
مكافئًا بحسب~$\beta$ للحد~$gk$. نعرّف الحدين
\[
\fn{diag}(x) = xx
\]
و
\[
l(x) = g(\fn{diag}(x))
\]
على مقتضى اصطلاحاتنا؛ فـ$l$ هو الحد~$\lambd[x][g(xx)]$.
ثم نأخذ~$k$ مساويًا للحد~$ll$، فنجد
\begin{align*}
k & = (\lambd[x][g(xx)])(\lambd[x][g(xx)]) \\
& \red  g((\lambd[x][g(xx)])(\lambd[x][g(xx)])) \\
& = gk.
\end{align*}
وعليه، إذا وضعنا
\[
Y = \lambd[g][((\lambd[x][g(xx)])(\lambd[x][g(xx)]))]
\]
كان~$Yg$ و$g(Yg)$ قابلين للاختزال إلى حد واحد، فينتج
$Yg \equiv_\beta g(Yg)$. وهذا هو «مؤلف كاري». وأما إذا وضعنا
\[
Y = (\lambd[xg][g(xxg)])(\lambd[xg][g(xxg)])
\]
فإن~$Yg$ يختزل إلى~$g(Yg)$ نفسه؛ وهذه نتيجة أقوى من مجرد التكافؤ.
وتسمى هذه الصيغة الأخيرة من~$Y$ «مؤلف تورنغ».

\end{document}

